\documentclass[11pt]{article}
\usepackage[margin=1in]{geometry}
\usepackage[T1]{fontenc}
\usepackage{amsmath,amssymb,mathtools}
\usepackage{booktabs,array,enumitem,caption}
\usepackage{xcolor,hyperref}
\sloppy
\hypersetup{colorlinks=true, linkcolor=blue!50!black, urlcolor=blue!50!black}
\newcommand{\code}[1]{\texttt{#1}}
\newcommand{\file}[1]{\texttt{#1}}
\setlist{nosep, leftmargin=1.4em}

\title{AutoBayes against brms: speed on ten models}
\author{Benchmark run of 17 October 2026, prepared for Andrew}
\date{17 October 2026}

\begin{document}
\maketitle

\noindent \textbf{The answer.} Yes, the speedup holds, and it holds for the Tier C and D models as well as the four
from the status report. Fitting each model from scratch (compiling and sampling), AutoBayes was 2.3 to 12.7 times faster
than brms, with a median of 5.9 times. Counting sampling alone, it was 1.2 to 7.9 times faster (median 3.9). Its
effective sample size per second of sampling was higher on nine of the ten models. On the same fits brms's
post-processing functions were 2 to 7 times slower than AutoBayes's. The two libraries' posteriors agree. The four
models that were also in the status report reproduced its numbers within about 10\% for brms; for AutoBayes they
reproduced within the same margin except B1, which was faster this time (5.3 s against 8.1 s).

\section{What was run}

Each model was fitted by both libraries with the same formula, family and explicit priors (no defaults), on the same
simulated data (\file{examples/brms\_parity/benchmark\_data.py}, fixed seeds, every numeric predictor standardised).
The sampler settings were also the same: four chains in parallel, 1000 warm-up and 1000 kept draws each, target
acceptance 0.9.
\begin{itemize}
\item \textbf{brms} 2.23.0 with rstan 2.32.7. Compilation (\code{brm(chains = 0)}) and sampling
  (\code{update(recompile = FALSE)}) were timed separately.
\item \textbf{AutoBayes} sampled with NumPyro's NUTS on the CPU (\code{chain\_method="parallel"}). Its time is the
  model build plus the fit, and the fit includes JAX's compilation. Its ``sampling'' time is the fit minus a one-draw
  fit of the same model, which estimates the compilation.
\item The machine was an Apple M1 Pro laptop in ordinary use, with desktop applications open and a Time Machine backup
  running when the run started. Each model was run once, brms first and then AutoBayes, never both at once.
\end{itemize}

\section{Time to a posterior}

\begin{table}[h]
\footnotesize
\setlength{\tabcolsep}{4pt}
\begin{center}\begin{tabular}{@{}llrrrrrrr@{}}
\toprule
& model & rows & \multicolumn{3}{c}{brms (s)} & AutoBayes (s) & \multicolumn{2}{c}{AutoBayes faster by} \\
\cmidrule(lr){4-6}\cmidrule(lr){8-9}
& & & compile & sample & total & total & total & sampling \\
\midrule
B1 & Bernoulli, \code{(1 + x | g)} & 2,000 & 36.2 & 5.8 & 42.0 & 5.3 & 7.9 & 5.4 \\
B2 & Gaussian, \code{(1 + x | g)} & 20,000 & 35.5 & 87.0 & 122.5 & 29.6 & 4.1 & 5.3 \\
B3 & cumulative, five categories & 1,500 & 27.0 & 33.4 & 60.4 & 7.1 & 8.5 & 7.9 \\
B4 & zero-inflated Poisson & 2,000 & 26.2 & 15.0 & 41.2 & 5.9 & 7.0 & 4.1 \\
B5 & Weibull & 2,000 & 27.1 & 6.4 & 33.5 & 6.0 & 5.6 & 3.3 \\
B6 & zero-inflated neg.\ binomial & 2,000 & 26.6 & 16.4 & 43.0 & 14.7 & 2.9 & 1.4 \\
B7 & sratio, five categories & 1,500 & 26.1 & 17.6 & 43.6 & 7.1 & 6.2 & 3.7 \\
B8 & Gaussian, horseshoe on 20 & 400 & 26.6 & 3.0 & 29.6 & 5.3 & 5.6 & 1.3 \\
B9 & beta-binomial, trials & 2,000 & 27.7 & 24.4 & 52.1 & 22.3 & 2.3 & 1.2 \\
B10 & Gaussian, \code{(1 | mm(g1, g2))} & 1,000 & 26.8 & 3.2 & 30.0 & 2.4 & 12.7 & 6.6 \\
\bottomrule
\end{tabular}\end{center}
\caption{\small Wall times, one run each. B1, B2, B3 and B5--B7 and B9 have a varying intercept (and B1, B2 a slope) on 30
to 200 groups.}\label{tab:time}
\end{table}

\noindent Most of the end-to-end difference on small models is brms's compilation. Turning a brms model into a C++
program and compiling it took 26 to 36 seconds per model; AutoBayes's JAX compilation took 2 to 4. On B10, for example,
brms spent 26.8 of its 30.0 seconds compiling and 3.2 sampling, while AutoBayes took 2.4 seconds in all. Sampling
itself was also faster in AutoBayes, but by less, and by very different amounts:
\begin{itemize}
\item \textbf{Large gains (5--8$\times$)} came on the models with many rows or an ordinal likelihood: B1--B3 and B10. On
  B2 (20,000 rows) sampling took 16.5 seconds against brms's 87.0.
\item \textbf{Small gains (1.2--1.4$\times$)} came on B6, B8 and B9. On B8 both samplers finish in about 2--3 seconds,
  so fixed overheads dominate. B6 and B9 evaluate log-gamma functions and their derivatives on every row (the negative
  binomial and the beta-binomial). That is a plausible cause of the small gain there, but I have not profiled it.
\end{itemize}

\section{Is the faster sampler sampling as well?}

Raw speed means little if the draws are worse. Table~\ref{tab:eff} reports the smallest bulk effective sample size over
the parameters both libraries report, computed by the same ArviZ estimator for both, and that number divided by the
sampling time.

\begin{table}[h]
\small
\begin{center}\begin{tabular}{@{}lrrrrrrr@{}}
\toprule
& \multicolumn{2}{c}{min.\ bulk ESS} & \multicolumn{2}{c}{ESS per second of sampling} & \multicolumn{2}{c}{divergent} & max $|z|$ \\
\cmidrule(lr){2-3}\cmidrule(lr){4-5}\cmidrule(lr){6-7}
& AutoBayes & brms & AutoBayes & brms & AutoBayes & brms & \\
\midrule
B1 & 1264 & 1017 & 1184 & 176 & 0 & 0 & 1.88 \\
B2 & 125 & 144 & 8 & 2 & 0 & 0 & 1.22 \\
B3 & 530 & 558 & 126 & 17 & 0 & 0 & 0.41 \\
B4 & 386 & 414 & 106 & 28 & 0 & 0 & 1.46 \\
B5 & 201 & 287 & 105 & 45 & 0 & 0 & 1.02 \\
B6 & 1079 & 1245 & 90 & 76 & 0 & 0 & 0.70 \\
B7 & 1087 & 1195 & 230 & 68 & 0 & 0 & 1.86 \\
B8 & 1916 & 2257 & 841 & 752 & 9 & 6 & 2.71 \\
B9 & 185 & 264 & 9 & 11 & 0 & 0 & 1.43 \\
B10 & 533 & 499 & 1093 & 154 & 0 & 0 & 0.82 \\
\bottomrule
\end{tabular}\end{center}
\caption{\small Sampling efficiency and agreement. $z$ is the difference of the two posterior means over their combined
Monte Carlo standard error, for each parameter both report; the largest over the parameters is shown.}\label{tab:eff}
\end{table}

\begin{itemize}
\item The ESS per second favours AutoBayes on nine models. Its smallest ESS is sometimes lower (B5: 201 against 287; B9:
  185 against 264), but it gets there faster. On B5, 201 effective draws in 1.9 seconds is about 105 a second, against
  brms's 287 in 6.4 seconds, about 45.
\item The exception is B9, the beta-binomial: 9 effective draws a second against brms's 11. On sampling efficiency
  alone, brms is ahead here. In every model the slowest parameter for both libraries is the population intercept (a
  threshold for the ordinal models, a coefficient on B8). That is probably the familiar slow mixing of the intercept
  against the mean of the varying intercepts, and it is the same in both libraries.
\item The posteriors agree: every parameter's mean differs by less than 2.8 combined Monte Carlo standard errors, and by
  less than 1.9 except on the horseshoe regression (2.71). That is the agreement expected from two independent samplers.
\item Both libraries hit divergent transitions only on the horseshoe regression (9 and 6 of 4,000). They are
  characteristic of that prior's funnel, not of either library.
\end{itemize}

\section{Post-processing}

Table~\ref{tab:post} times brms's functions on each fit against AutoBayes's: the median of three calls after a first one.

\begin{table}[h]
\footnotesize
\setlength{\tabcolsep}{3.5pt}
\begin{center}\begin{tabular}{@{}lccccccc@{}}
\toprule
& \code{posterior\_epred} & \code{log\_lik} & \code{loo} & \code{waic} & \code{bayes\_R2} & \code{hypothesis} & \code{cond.\_effects} \\
\midrule
B1 & 0.11 / 0.37 & 0.51 / 0.75 & 2.10 / 3.41 & 0.62 / 2.66 & 0.16 / 0.89 & 0.00 / 0.00 & 0.02 / 0.17 \\
B2 & 0.65 / 1.97 & 2.66 / 5.81 & 18.62 / 38.94 & 3.97 / 20.30 & 0.94 / 3.68 & 0.00 / 0.00 & 0.02 / 0.15 \\
B3 & 1.77 / 1.75 & 0.56 / 1.27 & 1.69 / 3.58 & 0.64 / 2.51 & 1.91 / 2.37 & \textit{ref.} / 0.00 & \textit{ref.} / 0.25 \\
B4 & 0.13 / 0.35 & 0.34 / 1.50 & 1.93 / 3.93 & 0.45 / 2.77 & 0.16 / 0.65 & 0.00 / 0.00 & 0.02 / 0.13 \\
B5 & 0.11 / 0.36 & 0.18 / 0.98 & 1.91 / 3.57 & 0.30 / 2.13 & 0.15 / 0.65 & 0.00 / 0.00 & 0.02 / 0.14 \\
B6 & 0.13 / 0.35 & 0.98 / 1.77 & 2.59 / 4.25 & 1.08 / 2.96 & 0.16 / 0.64 & 0.00 / 0.00 & 0.02 / 0.13 \\
B7 & 1.29 / 9.63 & 1.19 / 1.98 & 2.35 / 4.27 & 1.32 / 3.11 & 1.47 / 10.47 & \textit{ref.} / 0.00 & \textit{ref.} / 0.79 \\
B8 & 0.03 / 0.06 & 0.03 / 0.10 & 0.31 / 0.60 & 0.05 / 0.36 & 0.04 / 0.13 & 0.00 / 0.00 & 0.03 / 0.22 \\
B9 & 0.14 / 0.41 & 1.45 / \textit{n/a} & 3.06 / \textit{n/a} & 1.56 / \textit{n/a} & 0.18 / 0.70 & 0.00 / 0.00 & \textit{ref.} / 0.14 \\
B10 & 0.05 / 0.09 & 0.07 / 0.36 & 0.79 / 1.81 & 0.13 / 1.00 & 0.06 / 0.20 & 0.00 / 0.00 & 0.02 / 0.14 \\
\bottomrule
\end{tabular}\end{center}
\caption{\small Seconds, AutoBayes / brms. \textit{ref.}: AutoBayes refuses the function for this model, by design
(divergence D99: the ordinal families' thresholds, and the trials a beta-binomial grid would need). \textit{n/a}: brms
needs the R package \code{extraDistr} for the beta-binomial's \code{log\_lik}, and it is not installed here.}\label{tab:post}
\end{table}

\noindent The median ratios over the models where both ran were: \code{posterior\_epred} 2.8$\times$,
\code{log\_lik} 2.3$\times$, \code{loo} 2.0$\times$, \code{waic} 5.1$\times$, \code{bayes\_R2} 4.0$\times$ and
\code{conditional\_effects} 6.9$\times$. \code{hypothesis} takes a few milliseconds at most in both. The largest absolute
gap is PSIS-LOO on B2: 18.6 seconds against 38.9, for 20,000 observations and 4,000 draws.

\section{What these numbers do not show}

\begin{itemize}
\item \textbf{One run each.} The run was on a laptop in normal use, so a single number may be off by 10--30\%. The four
  models repeated from the status report came out within that range of their earlier times.
\item \textbf{Repeated fits.} brms compiles once per model and can refit new data without recompiling
  (\code{update(recompile = FALSE)}). For many fits of one model, its 26--36 seconds are paid once and only the
  ``sampling'' column matters.
\item \textbf{Other Stan setups.} cmdstanr, which runs a newer Stan and may sample faster than rstan, is not installed
  here and was not tested. Neither was brms's within-chain threading.
\item \textbf{GPUs.} The CPU only: no GPU result is claimed here.
\item \textbf{Scope.} Ten models on simulated data from each model's own family. A model that is hard to sample (a
  funnel, weak identification) is a different test.
\end{itemize}

\section{Reproducing}

\begin{verbatim}
python examples/brms_parity/benchmark_data.py DATA
Rscript examples/brms_parity/benchmark_brms.R DATA OUT
python examples/brms_parity/benchmark_autobayes.py DATA OUT
python examples/brms_parity/compare.py OUT
\end{verbatim}
The summary of this run is \file{examples/brms\_parity/results/comparison-2026-10-17.json}; the draws are not committed.

\end{document}
